\scanpage{21}
% archive sheet number: 20

1) $\forall\,0\leq i\leq2$, $\tau_i^C$ est $\ne 1$ en chaque fibre $\left(\Rightarrow
\begin{array}{ccc}h_0,&h_1,&h_2\\ H_0,&H_1,&H_2\\\multicolumn{3}{c}{\text{dans } C}\end{array}\right]$

2) $h_0\ h_1\ h_2$ lin. indép. en chaque fibre

\deleted{(à \unsure{expliciter} \inserted{\unsure{plus intrinsèque}} en termes de $\unsure{\varphi_G^Y}$)}

$\Longleftrightarrow$ \deleted{$\forall$} $\forall$ fibre \deleted{\unsure{générique}} $Y_{\bar s}$, il \uline{n'existe pas} de diviseur $\Delta$ de $Y_{\bar s}$ de degré $2$, invariant par $\varphi_G^Y$, i.e. pas de restriction à une structure soit de droite affine, soit de \deleted{droite affine} torseur sous $\mathrm G_m$ \ldots\ (avec $G$ agissant sur $\mathrm G_m$ \ldots)

% slanted
\margin{C'est sans doute ce qui exclut le cas ``\unsure{diédral|dièdral}''}

3) $\tau_0^C\tau_1^C\overset{\mathrm{dfn}}=\rho_f$ et $\tau_1^C\tau_2^C\overset{\mathrm{dfn}}=\rho_s$ sont tels que
\[
\rho_f^2\ne\mathrm{id}_{\unsure{Y|C}},\quad \rho_s^2\ne\mathrm{id}_Y\quad\text{sur chaque fibre.}
\]

(cela exprime la non commutation de $\tau_0,\tau_1$, et de $\tau_1,\tau_2$, sur chaque fibre
